% --- Archon physics formalization source begin ---
% archon:physics
% archon:covers PhyXMiniProblems/problem_phyx_mini_0805.lean
% archon:source-report reports/phyx_mini/problem_phyx_mini_0805.source.json

\chapter{Physics problem phyx\_mini\_0805}
\label{ch:PhyXMiniProblems_problem_phyx_mini_0805}

\paragraph{Problem source.}
\#\# Physical scenario

Figure shows an air-track glider with mass \textbackslash{}( m\_1 \textbackslash{}) moving on a level, frictionless air track in the physics lab. The glider is connected to a lab weight with mass \textbackslash{}( m\_2 \textbackslash{}) by a light, flexible, nonstretching string that passes over a stationary, frictionless pulley.

\#\# Question

Find the tension in the string.

\#\# Answer choices

- A: A: \textbackslash{}( \textbackslash{}frac\{2 m\_1 m\_2\}\{m\_1 + m\_2\} g \textbackslash{})

- B: B: \textbackslash{}( \textbackslash{}frac\{m\_1 m\_2\}\{m\_1 + m\_2\} g \textbackslash{})

- C: C: \textbackslash{}( \textbackslash{}frac\{m\_1 m\_2\}\{2 m\_1 + m\_2\} g \textbackslash{})

- D: D: \textbackslash{}( \textbackslash{}frac\{m\_1 m\_2\}\{m\_1 + 2 m\_2\} g \textbackslash{})

\#\# Figure caption (auxiliary; use the image as primary evidence)

The image depicts a physics apparatus often used to demonstrate concepts of mechanics, specifically involving pulleys and inclined planes. Here's a detailed breakdown:

- There is an inclined plane with a dotted surface, indicating friction or texture for traction.
- An object labeled \textbackslash{}( m\_1 \textbackslash{}) is placed on the inclined plane.
- A pulley system is shown at the top right corner of the plane.
- A string or cable is connected to \textbackslash{}( m\_1 \textbackslash{}) and goes over the pulley.
- The string is attached to another object labeled \textbackslash{}( m\_2 \textbackslash{}), which is hanging vertically off the edge of the pulley, showing a counterweight setup.
- The setup is commonly used to study the forces and motion in systems with friction and tension.

\#\# Dataset metadata

Dynamics; Physical Model Grounding Reasoning, Multi-Formula Reasoning

\paragraph{Recorded answer/context.}
B: B: \textbackslash{}( \textbackslash{}frac\{m\_1 m\_2\}\{m\_1 + m\_2\} g \textbackslash{})

\paragraph{Figure/image path.}
/root/proposal\_for\_physic/science-mango/phyx\_mini\_run/phyx\_data/test\_image/805.png

\paragraph{Formalization target.}
create a compiling Lean file with sorry bodies at `PhyXMiniProblems/problem\_phyx\_mini\_0805.lean`. The Lean declarations must preserve the physical quantities, dimensions or dimensional roles, figure labels, governing-law hypotheses, and final relation expressed by this problem.
Use Mathlib/Physlib names found through LeanExplore where available. If a physics API is missing, introduce faithful local abstractions rather than scalar placeholder aliases.

\begin{theorem}[Physics formalization target]
\label{thm:physics:phyx_mini_0805:target}
\lean{PhyXMiniProblems.ProblemPhyXMini0805.problem_phyx_mini_0805}
\uses{def:physics:phyx-mini-0805:phyxminiproblems-problemphyxmini0805-massinkilograms, def:physics:phyx-mini-0805:phyxminiproblems-problemphyxmini0805-accelerationinmeterspersecondsquared, def:physics:phyx-mini-0805:phyxminiproblems-problemphyxmini0805-tensioninnewtons, def:physics:phyx-mini-0805:phyxminiproblems-problemphyxmini0805-airtrackpulleysetup, def:physics:phyx-mini-0805:phyxminiproblems-problemphyxmini0805-matcheswrittenscenario, def:physics:phyx-mini-0805:phyxminiproblems-problemphyxmini0805-matchessuppliedfigure, def:physics:phyx-mini-0805:phyxminiproblems-problemphyxmini0805-haspositivephysicalparameters, def:physics:phyx-mini-0805:phyxminiproblems-problemphyxmini0805-satisfiesidealairtrackpulleylaws, def:physics:phyx-mini-0805:phyxminiproblems-problemphyxmini0805-recordeddatasetanswer, def:physics:phyx-mini-0805:phyxminiproblems-problemphyxmini0805-isuniquematchinganswer}
The assigned autoformalize agent should translate this physics problem into Lean declarations in the covered file, with theorem and lemma proof bodies written as `by sorry`.
\end{theorem}
\begin{proof}
This is an autoformalization task, not a proof task. Produce statements that can later be proved by the physics prover without weakening the physical model.
\end{proof}

% --- Archon declaration topology begin ---
\section{Declaration topology}
\begin{definition}[acceleration Dimension]
  \label{def:physics:phyx-mini-0805:phyxminiproblems-problemphyxmini0805-accelerationdimension}
  \lean{PhyXMiniProblems.ProblemPhyXMini0805.accelerationDimension}
  The physical dimension L T⁻² of an acceleration magnitude
\end{definition}
\begin{proof}
  This is the declared model definition of acceleration Dimension.
\end{proof}
\begin{definition}[force Dimension]
  \label{def:physics:phyx-mini-0805:phyxminiproblems-problemphyxmini0805-forcedimension}
  \lean{PhyXMiniProblems.ProblemPhyXMini0805.forceDimension}
  The physical dimension M L T⁻² of a force magnitude
\end{definition}
\begin{proof}
  This is the declared model definition of force Dimension.
\end{proof}
\begin{definition}[Mass Quantity]
  \label{def:physics:phyx-mini-0805:phyxminiproblems-problemphyxmini0805-massquantity}
  \lean{PhyXMiniProblems.ProblemPhyXMini0805.MassQuantity}
  A nonnegative, unit-independent physical mass
\end{definition}
\begin{proof}
  This is the declared model definition of Mass Quantity.
\end{proof}
\begin{definition}[Acceleration Quantity]
  \label{def:physics:phyx-mini-0805:phyxminiproblems-problemphyxmini0805-accelerationquantity}
  \lean{PhyXMiniProblems.ProblemPhyXMini0805.AccelerationQuantity}
  \uses{def:physics:phyx-mini-0805:phyxminiproblems-problemphyxmini0805-accelerationdimension}
  A nonnegative, unit-independent physical acceleration magnitude
\end{definition}
\begin{proof}
  This is the declared model definition of Acceleration Quantity.
\end{proof}
\begin{definition}[Tension Quantity]
  \label{def:physics:phyx-mini-0805:phyxminiproblems-problemphyxmini0805-tensionquantity}
  \lean{PhyXMiniProblems.ProblemPhyXMini0805.TensionQuantity}
  \uses{def:physics:phyx-mini-0805:phyxminiproblems-problemphyxmini0805-forcedimension}
  A nonnegative, unit-independent force magnitude, used for string tension
\end{definition}
\begin{proof}
  This is the declared model definition of Tension Quantity.
\end{proof}
\begin{definition}[mass Readout]
  \label{def:physics:phyx-mini-0805:phyxminiproblems-problemphyxmini0805-massreadout}
  \lean{PhyXMiniProblems.ProblemPhyXMini0805.massReadout}
  \uses{def:physics:phyx-mini-0805:phyxminiproblems-problemphyxmini0805-massquantity}
  Read a mass in the mass unit selected by a coherent unit system
\end{definition}
\begin{proof}
  This is the declared model definition of mass Readout.
\end{proof}
\begin{definition}[acceleration Readout]
  \label{def:physics:phyx-mini-0805:phyxminiproblems-problemphyxmini0805-accelerationreadout}
  \lean{PhyXMiniProblems.ProblemPhyXMini0805.accelerationReadout}
  \uses{def:physics:phyx-mini-0805:phyxminiproblems-problemphyxmini0805-accelerationquantity}
  Read an acceleration in the coherent unit selected by a unit system
\end{definition}
\begin{proof}
  This is the declared model definition of acceleration Readout.
\end{proof}
\begin{definition}[tension Readout]
  \label{def:physics:phyx-mini-0805:phyxminiproblems-problemphyxmini0805-tensionreadout}
  \lean{PhyXMiniProblems.ProblemPhyXMini0805.tensionReadout}
  \uses{def:physics:phyx-mini-0805:phyxminiproblems-problemphyxmini0805-tensionquantity}
  Read a tension in the coherent force unit selected by a unit system
\end{definition}
\begin{proof}
  This is the declared model definition of tension Readout.
\end{proof}
\begin{definition}[mass In Kilograms]
  \label{def:physics:phyx-mini-0805:phyxminiproblems-problemphyxmini0805-massinkilograms}
  \lean{PhyXMiniProblems.ProblemPhyXMini0805.massInKilograms}
  \uses{def:physics:phyx-mini-0805:phyxminiproblems-problemphyxmini0805-massquantity, def:physics:phyx-mini-0805:phyxminiproblems-problemphyxmini0805-massreadout}
  Kilogram readout of a physical mass
\end{definition}
\begin{proof}
  This is the declared model definition of mass In Kilograms.
\end{proof}
\begin{definition}[acceleration In Meters Per Second Squared]
  \label{def:physics:phyx-mini-0805:phyxminiproblems-problemphyxmini0805-accelerationinmeterspersecondsquared}
  \lean{PhyXMiniProblems.ProblemPhyXMini0805.accelerationInMetersPerSecondSquared}
  \uses{def:physics:phyx-mini-0805:phyxminiproblems-problemphyxmini0805-accelerationquantity, def:physics:phyx-mini-0805:phyxminiproblems-problemphyxmini0805-accelerationreadout}
  Metres-per-second-squared readout of an acceleration magnitude
\end{definition}
\begin{proof}
  This is the declared model definition of acceleration In Meters Per Second Squared.
\end{proof}
\begin{definition}[tension In Newtons]
  \label{def:physics:phyx-mini-0805:phyxminiproblems-problemphyxmini0805-tensioninnewtons}
  \lean{PhyXMiniProblems.ProblemPhyXMini0805.tensionInNewtons}
  \uses{def:physics:phyx-mini-0805:phyxminiproblems-problemphyxmini0805-tensionquantity, def:physics:phyx-mini-0805:phyxminiproblems-problemphyxmini0805-tensionreadout}
  Newton readout of a string-tension magnitude
\end{definition}
\begin{proof}
  This is the declared model definition of tension In Newtons.
\end{proof}
\begin{definition}[Body Label]
  \label{def:physics:phyx-mini-0805:phyxminiproblems-problemphyxmini0805-bodylabel}
  \lean{PhyXMiniProblems.ProblemPhyXMini0805.BodyLabel}
  The two bodies labelled in the problem and the supplied raster
\end{definition}
\begin{proof}
  This is the declared model definition of Body Label.
\end{proof}
\begin{definition}[String Segment]
  \label{def:physics:phyx-mini-0805:phyxminiproblems-problemphyxmini0805-stringsegment}
  \lean{PhyXMiniProblems.ProblemPhyXMini0805.StringSegment}
  The two straight portions of the pictured string
\end{definition}
\begin{proof}
  This is the declared model definition of String Segment.
\end{proof}
\begin{definition}[Figure Object]
  \label{def:physics:phyx-mini-0805:phyxminiproblems-problemphyxmini0805-figureobject}
  \lean{PhyXMiniProblems.ProblemPhyXMini0805.FigureObject}
  Individually recognizable objects visible in the primary raster
\end{definition}
\begin{proof}
  This is the declared model definition of Figure Object.
\end{proof}
\begin{definition}[Figure Text Label]
  \label{def:physics:phyx-mini-0805:phyxminiproblems-problemphyxmini0805-figuretextlabel}
  \lean{PhyXMiniProblems.ProblemPhyXMini0805.FigureTextLabel}
  Literal symbolic mass labels visible in the primary raster
\end{definition}
\begin{proof}
  This is the declared model definition of Figure Text Label.
\end{proof}
\begin{definition}[Body Location]
  \label{def:physics:phyx-mini-0805:phyxminiproblems-problemphyxmini0805-bodylocation}
  \lean{PhyXMiniProblems.ProblemPhyXMini0805.BodyLocation}
  Locations of the two labelled bodies in the apparatus
\end{definition}
\begin{proof}
  This is the declared model definition of Body Location.
\end{proof}
\begin{definition}[Segment Orientation]
  \label{def:physics:phyx-mini-0805:phyxminiproblems-problemphyxmini0805-segmentorientation}
  \lean{PhyXMiniProblems.ProblemPhyXMini0805.SegmentOrientation}
  Geometric orientation of a straight string segment
\end{definition}
\begin{proof}
  This is the declared model definition of Segment Orientation.
\end{proof}
\begin{definition}[String Route]
  \label{def:physics:phyx-mini-0805:phyxminiproblems-problemphyxmini0805-stringroute}
  \lean{PhyXMiniProblems.ProblemPhyXMini0805.StringRoute}
  The complete route of the single string shown in the figure
\end{definition}
\begin{proof}
  This is the declared model definition of String Route.
\end{proof}
\begin{definition}[Track Orientation]
  \label{def:physics:phyx-mini-0805:phyxminiproblems-problemphyxmini0805-trackorientation}
  \lean{PhyXMiniProblems.ProblemPhyXMini0805.TrackOrientation}
  Orientation of the air track stated in the problem
\end{definition}
\begin{proof}
  This is the declared model definition of Track Orientation.
\end{proof}
\begin{definition}[Track Resistance Model]
  \label{def:physics:phyx-mini-0805:phyxminiproblems-problemphyxmini0805-trackresistancemodel}
  \lean{PhyXMiniProblems.ProblemPhyXMini0805.TrackResistanceModel}
  Resistance model for glider motion along the track
\end{definition}
\begin{proof}
  This is the declared model definition of Track Resistance Model.
\end{proof}
\begin{definition}[String Mass Model]
  \label{def:physics:phyx-mini-0805:phyxminiproblems-problemphyxmini0805-stringmassmodel}
  \lean{PhyXMiniProblems.ProblemPhyXMini0805.StringMassModel}
  Mass idealization for the connecting string
\end{definition}
\begin{proof}
  This is the declared model definition of String Mass Model.
\end{proof}
\begin{definition}[String Flexibility Model]
  \label{def:physics:phyx-mini-0805:phyxminiproblems-problemphyxmini0805-stringflexibilitymodel}
  \lean{PhyXMiniProblems.ProblemPhyXMini0805.StringFlexibilityModel}
  Flexibility idealization for the connecting string
\end{definition}
\begin{proof}
  This is the declared model definition of String Flexibility Model.
\end{proof}
\begin{definition}[String Extensibility Model]
  \label{def:physics:phyx-mini-0805:phyxminiproblems-problemphyxmini0805-stringextensibilitymodel}
  \lean{PhyXMiniProblems.ProblemPhyXMini0805.StringExtensibilityModel}
  Extensibility idealization for the connecting string
\end{definition}
\begin{proof}
  This is the declared model definition of String Extensibility Model.
\end{proof}
\begin{definition}[Pulley Motion Model]
  \label{def:physics:phyx-mini-0805:phyxminiproblems-problemphyxmini0805-pulleymotionmodel}
  \lean{PhyXMiniProblems.ProblemPhyXMini0805.PulleyMotionModel}
  Motion state of the pulley support
\end{definition}
\begin{proof}
  This is the declared model definition of Pulley Motion Model.
\end{proof}
\begin{definition}[Pulley Resistance Model]
  \label{def:physics:phyx-mini-0805:phyxminiproblems-problemphyxmini0805-pulleyresistancemodel}
  \lean{PhyXMiniProblems.ProblemPhyXMini0805.PulleyResistanceModel}
  Axle/contact resistance model for the pulley
\end{definition}
\begin{proof}
  This is the declared model definition of Pulley Resistance Model.
\end{proof}
\begin{definition}[Positive Axis Direction]
  \label{def:physics:phyx-mini-0805:phyxminiproblems-problemphyxmini0805-positiveaxisdirection}
  \lean{PhyXMiniProblems.ProblemPhyXMini0805.PositiveAxisDirection}
  Positive coordinate direction used in each body's scalar force equation
\end{definition}
\begin{proof}
  This is the declared model definition of Positive Axis Direction.
\end{proof}
\begin{definition}[Supplied Air Track Figure]
  \label{def:physics:phyx-mini-0805:phyxminiproblems-problemphyxmini0805-suppliedairtrackfigure}
  \lean{PhyXMiniProblems.ProblemPhyXMini0805.SuppliedAirTrackFigure}
  \uses{def:physics:phyx-mini-0805:phyxminiproblems-problemphyxmini0805-bodylabel, def:physics:phyx-mini-0805:phyxminiproblems-problemphyxmini0805-stringsegment, def:physics:phyx-mini-0805:phyxminiproblems-problemphyxmini0805-figureobject, def:physics:phyx-mini-0805:phyxminiproblems-problemphyxmini0805-figuretextlabel, def:physics:phyx-mini-0805:phyxminiproblems-problemphyxmini0805-bodylocation, def:physics:phyx-mini-0805:phyxminiproblems-problemphyxmini0805-segmentorientation, def:physics:phyx-mini-0805:phyxminiproblems-problemphyxmini0805-stringroute}
  Typed transcription of image 805.png. It records the horizontal air-track geometry actually visible in the raster; the auxiliary generated caption's claim that the track is inclined is deliberately not encoded
\end{definition}
\begin{proof}
  This is the declared model definition of Supplied Air Track Figure.
\end{proof}
\begin{definition}[Air Track Pulley Setup]
  \label{def:physics:phyx-mini-0805:phyxminiproblems-problemphyxmini0805-airtrackpulleysetup}
  \lean{PhyXMiniProblems.ProblemPhyXMini0805.AirTrackPulleySetup}
  \uses{def:physics:phyx-mini-0805:phyxminiproblems-problemphyxmini0805-massquantity, def:physics:phyx-mini-0805:phyxminiproblems-problemphyxmini0805-accelerationquantity, def:physics:phyx-mini-0805:phyxminiproblems-problemphyxmini0805-tensionquantity, def:physics:phyx-mini-0805:phyxminiproblems-problemphyxmini0805-bodylabel, def:physics:phyx-mini-0805:phyxminiproblems-problemphyxmini0805-stringsegment, def:physics:phyx-mini-0805:phyxminiproblems-problemphyxmini0805-trackorientation, def:physics:phyx-mini-0805:phyxminiproblems-problemphyxmini0805-trackresistancemodel, def:physics:phyx-mini-0805:phyxminiproblems-problemphyxmini0805-stringmassmodel, def:physics:phyx-mini-0805:phyxminiproblems-problemphyxmini0805-stringflexibilitymodel, def:physics:phyx-mini-0805:phyxminiproblems-problemphyxmini0805-stringextensibilitymodel, def:physics:phyx-mini-0805:phyxminiproblems-problemphyxmini0805-pulleymotionmodel, def:physics:phyx-mini-0805:phyxminiproblems-problemphyxmini0805-pulleyresistancemodel, def:physics:phyx-mini-0805:phyxminiproblems-problemphyxmini0805-positiveaxisdirection, def:physics:phyx-mini-0805:phyxminiproblems-problemphyxmini0805-suppliedairtrackfigure}
  Independent physical quantities and apparatus properties. The tension and accelerations are fields, rather than definitions manufactured from the recorded answer formula
\end{definition}
\begin{proof}
  This is the declared model definition of Air Track Pulley Setup.
\end{proof}
\begin{definition}[Matches Written Scenario]
  \label{def:physics:phyx-mini-0805:phyxminiproblems-problemphyxmini0805-matcheswrittenscenario}
  \lean{PhyXMiniProblems.ProblemPhyXMini0805.MatchesWrittenScenario}
  \uses{def:physics:phyx-mini-0805:phyxminiproblems-problemphyxmini0805-airtrackpulleysetup}
  Qualitative idealizations stated in the problem prose
\end{definition}
\begin{proof}
  This is the declared model definition of Matches Written Scenario.
\end{proof}
\begin{definition}[Matches Supplied Figure]
  \label{def:physics:phyx-mini-0805:phyxminiproblems-problemphyxmini0805-matchessuppliedfigure}
  \lean{PhyXMiniProblems.ProblemPhyXMini0805.MatchesSuppliedFigure}
  \uses{def:physics:phyx-mini-0805:phyxminiproblems-problemphyxmini0805-airtrackpulleysetup}
  Geometry and symbolic labels read directly from the primary raster
\end{definition}
\begin{proof}
  This is the declared model definition of Matches Supplied Figure.
\end{proof}
\begin{definition}[Has Positive Physical Parameters]
  \label{def:physics:phyx-mini-0805:phyxminiproblems-problemphyxmini0805-haspositivephysicalparameters}
  \lean{PhyXMiniProblems.ProblemPhyXMini0805.HasPositivePhysicalParameters}
  \uses{def:physics:phyx-mini-0805:phyxminiproblems-problemphyxmini0805-massinkilograms, def:physics:phyx-mini-0805:phyxminiproblems-problemphyxmini0805-accelerationinmeterspersecondsquared, def:physics:phyx-mini-0805:phyxminiproblems-problemphyxmini0805-airtrackpulleysetup}
  Positivity assumptions selecting the nondegenerate physical regime
\end{definition}
\begin{proof}
  This is the declared model definition of Has Positive Physical Parameters.
\end{proof}
\begin{definition}[Satisfies Ideal Air Track Pulley Laws]
  \label{def:physics:phyx-mini-0805:phyxminiproblems-problemphyxmini0805-satisfiesidealairtrackpulleylaws}
  \lean{PhyXMiniProblems.ProblemPhyXMini0805.SatisfiesIdealAirTrackPulleyLaws}
  \uses{def:physics:phyx-mini-0805:phyxminiproblems-problemphyxmini0805-massreadout, def:physics:phyx-mini-0805:phyxminiproblems-problemphyxmini0805-accelerationreadout, def:physics:phyx-mini-0805:phyxminiproblems-problemphyxmini0805-tensionreadout, def:physics:phyx-mini-0805:phyxminiproblems-problemphyxmini0805-airtrackpulleysetup}
  The governing ideal-string, ideal-pulley, and Newton-law relations. The first two laws express the constraints supplied by a single light nonstretching string over a frictionless fixed pulley. The final two equations are Newton's second law along the chosen positive axes: T = m₁ a₁ for the frictionless glider and m₂ g - T = m₂ a₂ for the descending hanging mass. These are generic relations in every coherent unit system. None contains the solved tension formula asked for in the current question
\end{definition}
\begin{proof}
  This is the declared model definition of Satisfies Ideal Air Track Pulley Laws.
\end{proof}
\begin{definition}[Answer Choice]
  \label{def:physics:phyx-mini-0805:phyxminiproblems-problemphyxmini0805-answerchoice}
  \lean{PhyXMiniProblems.ProblemPhyXMini0805.AnswerChoice}
  Labels of the four symbolic alternatives printed with the problem
\end{definition}
\begin{proof}
  This is the declared model definition of Answer Choice.
\end{proof}
\begin{definition}[displayed Tension]
  \label{def:physics:phyx-mini-0805:phyxminiproblems-problemphyxmini0805-answerchoice-displayedtension}
  \lean{PhyXMiniProblems.ProblemPhyXMini0805.AnswerChoice.displayedTension}
  \uses{def:physics:phyx-mini-0805:phyxminiproblems-problemphyxmini0805-answerchoice}
  Scalar tension formula printed beside an answer label, evaluated from coherent readouts m₁, m₂, and g
\end{definition}
\begin{proof}
  This is the declared model definition of displayed Tension.
\end{proof}
\begin{definition}[recorded Dataset Answer]
  \label{def:physics:phyx-mini-0805:phyxminiproblems-problemphyxmini0805-recordeddatasetanswer}
  \lean{PhyXMiniProblems.ProblemPhyXMini0805.recordedDatasetAnswer}
  \uses{def:physics:phyx-mini-0805:phyxminiproblems-problemphyxmini0805-answerchoice}
  Dataset answer metadata, deliberately not used as a theorem premise
\end{definition}
\begin{proof}
  This is the declared model definition of recorded Dataset Answer.
\end{proof}
\begin{definition}[Matches Answer Choice]
  \label{def:physics:phyx-mini-0805:phyxminiproblems-problemphyxmini0805-matchesanswerchoice}
  \lean{PhyXMiniProblems.ProblemPhyXMini0805.MatchesAnswerChoice}
  \uses{def:physics:phyx-mini-0805:phyxminiproblems-problemphyxmini0805-massinkilograms, def:physics:phyx-mini-0805:phyxminiproblems-problemphyxmini0805-accelerationinmeterspersecondsquared, def:physics:phyx-mini-0805:phyxminiproblems-problemphyxmini0805-tensioninnewtons, def:physics:phyx-mini-0805:phyxminiproblems-problemphyxmini0805-airtrackpulleysetup, def:physics:phyx-mini-0805:phyxminiproblems-problemphyxmini0805-answerchoice}
  Every segment tension agrees with the formula printed for a choice
\end{definition}
\begin{proof}
  This is the declared model definition of Matches Answer Choice.
\end{proof}
\begin{definition}[Is Unique Matching Answer]
  \label{def:physics:phyx-mini-0805:phyxminiproblems-problemphyxmini0805-isuniquematchinganswer}
  \lean{PhyXMiniProblems.ProblemPhyXMini0805.IsUniqueMatchingAnswer}
  \uses{def:physics:phyx-mini-0805:phyxminiproblems-problemphyxmini0805-airtrackpulleysetup, def:physics:phyx-mini-0805:phyxminiproblems-problemphyxmini0805-answerchoice, def:physics:phyx-mini-0805:phyxminiproblems-problemphyxmini0805-matchesanswerchoice}
  A choice is the unique printed formula agreeing with the string tension
\end{definition}
\begin{proof}
  This is the declared model definition of Is Unique Matching Answer.
\end{proof}
% --- Archon declaration topology end ---
% --- Archon physics formalization source end ---
